\section{Algoritmo para elegir las monedas}

En este trabajo de ejemplo realizaremos el análisis teórico y empírico de diferentes algoritmos que sirven para resolver el mismo problema: obtener el orden en que se deben elegir las monedas para maximizar la ganancia obtenida de un arreglo de $n$ monedas. 

\subsection{Algoritmo para encontrar el optimo}

El algoritmo plantea una matriz de $n * n$, siendo $n$ la cantidad de monedas, donde para cada coordenada $i$, $j$ de la matriz, encontramos una solucion para un problema de dimensiones $i$, $j$.
Luego se llena las diagonales, donde las corrdenads $i$ y $j$ son iguales, $i = j$, y la diagonal donde la coordenada $j$ es uno mas que la $i$, $i = j + 1$, con los valores de los dos casos base, siendo estos:
\begin{itemize}
    \item tener una sola mondea, donde el optimo es la moneda en si.
    \item tener dos monedas, donde el optimo es la mas grande entre la dos.
\end{itemize}
Luego recorre la parte superior de la matriz (ya que los punteros $i$, $j$ que 
apuntan al principio y final del array de monedas respectivametne, nunca se
intercambian, lo que es igual, se cummple que $i <= j$), llenando cada lugar de la matriz con la solucion del problema de tamaño $i$, $j$.
Es asi que una vez completamos la matriz, encontramos la solucion optima para nuestro problema original en la coordenada $i = 0$ y $j = n$.

Aqui el codigo usado para generar la matriz con los valores optimos para todos los problemas de dimensiones $i$, $j$.

\begin{lstlisting}[language=Python]
def elegir(monedas):
    opt = [[0] * len(monedas) for i in range(len(monedas))]
    
    # Casos base con una moneda
    for i in range(0, len(monedas)):
        opt[i][i] = monedas[i]
        
    # Casos base con 2 monedas
    for i in range(0, len(monedas) - 1):
        j = i + 1
        opt[i][j] = max(monedas[i], monedas[j])
    
    for i in range(len(opt) - 1, -1, -1):
        for j in range(i + 2, len(opt)):
            # Sophia toma la primera
            # Achico mi rango segun lo que eliga Mateo
            opt_i = opt[i + 2][j] if monedas[i + 1] > monedas[j] else opt[i + 1][j - 1]
            
            # Sophia toma la ultima
            # Achico mi rango segun lo que eliga Mateo
            opt_j = opt[i + 1][j - 1] if monedas[i] > monedas[j - 1] else opt[i][j - 2]
            
            opt[i][j] = max(monedas[i] + opt_i, monedas[j] + opt_j)
            
    maxima_ganancia = opt[0][len(opt) - 1]
    sol = solucion(opt, monedas, [], 0, len(opt) - 1)
    return sol
\end{lstlisting}

La forma de obtener nuesta matriz optima es la siguente.
Se plantea que para cada problema de dimensiones $i$, $j$, donde queremos encontrar que moneda elegir para mazimizar la ganancia, la moneda en la posicion $i$ (la primera), y la moneda $j$ (la ultima), se puden analizar las siguentes opciones.

\begin{enumerate}
    \item Sophia toma la primera moneda, la moneda $i$ y Mateo elige la mayor entre las monedas $i + 1$ y $j$.

    \begin{enumerate}
        \item Si mateo eligie la moneda $i + 1$ el rango $(i, j)$ se achica a $(i + 2, j)$, por lo que nuestra solucion sera el valor de la moneda $i$ mas la solucion del rango $(i + 2, j)$.

        $opt(i, j) = moneda_i + opt(i + 2, j)$

        \item Si mateo eligie la moneda j el rango $(i, j)$ se achica a $(i + 1, j - 1)$, por lo que nuestra solucion sera el valor de la moneda $i$ mas la solucion del rango $(i + 1, j - j)$.

        $opt(i, j) = moneda_i + opt(i + 1, j - 1)$

    \end{enumerate}

    \item Sophia toma ahora la ultima moneda, la moneda $j$ y Mateo elige la mayor entre las monedas $i$ y $j - 1$.

    \begin{enumerate}
        \item Si mateo eligio la moneda $i$ el rango $(i, j)$ se achica a $(i + 1, j - 1)$, por lo que nuestra solucion sera el valor de la moneda j mas la solucion del rango $(i + 1, j - 1)$.

        $opt(i, j) = moneda_j + opt(i + 1, j - 1)$

        \item Si mateo eligio la moneda $j$ el rango $(i, j)$ se achica a $(i, j - 2)$, por lo que nuestra solucion sera el valor de la moneda j mas la solucion del rango $(i, j - 2)$.

        $opt(i, j) = moneda_i + opt(i, j - 2)$
    \end{enumerate}
\end{enumerate}

Luego de plantear las cuatro posibles situaciones, podemos ver que, 
sin importar lo que elija sophia, mateo siempre elije la moneda mas grande
por lo que en ambos casos en simplemente realizar una comparacion para ver cual 
es mas grande, segun la situacion de monedas. Esta comparacion determina el rango
de mis soluciones parciales. Luego obtengo mis soluciones parciales para los rangos 
respectivos segun la eleccion de mateo, me quedo con aquella solucion, que maximize mi 
ganancia. 
Entonces Siendo $opt_i$, la solucion para un rango $(i_i, j_i)$, asumiendo que sophia toma la $moneda i$, y, siendo $opt_j$, la solucion para un rango $(i_j, j_j)$, asumiendo que sophia toma la moneda $j$, el optimo para el rango $(i, j)$ queda de la siguente manera.

\[
\text{opt}(i, j) = \max \left( \text{moneda}_i + \text{opt}_i, \; \text{moneda}_j + \text{opt}_j \right)
\]

donde:
\[
\text{moneda}_i + \text{opt}_i = 
\begin{cases} 
\text{monedas}[i] + \text{opt}(i + 2, j) & \text{si } \text{monedas}[i + 1] > \text{monedas}[j] \\[1ex]
\text{monedas}[i] + \text{opt}(i + 1, j - 1) & \text{sino}
\end{cases}
\]

\[
\text{moneda}_j + \text{opt}_j = 
\begin{cases} 
\text{monedas}[j] + \text{opt}(i + 1, j - 1) & \text{si } \text{monedas}[i] > \text{monedas}[j - 1] \\[1ex]
\text{monedas}[j] + \text{opt}(i, j - 2) & \text{sino}
\end{cases}
\]

\subsection{Algoritmo para obtener la reconstruccion de la solucion}

Partiendo de la matriz optima, desarrollamos un algoritmo que brinde la secuencia de monedas que se deberian elegir para obtener la ganancia maxima, mediante llamadas recursivas.

\begin{enumerate}
    \item Casos base.
    \begin{itemize}
        \item $i < j$ El rango de monedas es invalido, finalizo la recursion y devuelvo la solucion.
        \item $i = j$ Solo queda una moneda por elejir, se agrega a la solucion y se devuelve el conjunto.
    \end{itemize}

    \item Evaluacion de alternativas (esta es la parte donde analizamos que sub problema conforma la solucion del problema actual y se simula la eleccion de monedas).
    \begin{itemize}
        \item Se calcula $opt_i$ y $opt_j$ de la misma manera que en el algoritmo anterior, siendo los optimos si se elije la moneda $i$ o la moneda $j$, y en concecuencia que moneda elijiria Mateo (la mayor), y como estas decisiones afectan las dimensiones de nuestro problema.
        \item Analizamos que optimo nos da la mayor ganancia si le sumamos la mondea correspondiente (porque es lo que el algoritmo hubiera hecho), y en base al mayor agremos el indice $i$ o el $j$ a la solucion, respectivametne.
        \item Por ultimo vemos que moneda elije Mateo, cual es la mayor, y llamamos a la funcion recursivamente con el nuevo rango $i$, $j$ conseguido luego de analizar las decisones que tomarian ambos bandos.
    \end{itemize}
\end{enumerate}

\begin{lstlisting}[language=Python]
def solucion(opt, monedas, s, i, j):
if i > j:
    return s

if i == j:
    s.append(i)
    return s

opt_i = opt[i + 2][j] if monedas[i + 1] > monedas[j] else opt[i + 1][j - 1]
opt_j = opt[i + 1][j - 1] if monedas[i] > monedas[j - 1] else opt[i][j - 2]

if monedas[i] + opt_i > monedas[j] + opt_j:
    s.append(i)
    
    if monedas[i + 1] > monedas[j]:
        return solucion(opt, monedas, s, i + 2, j)
    else:    
        return solucion(opt, monedas, s, i + 1, j - 1)
else:
    s.append(j)
    
    if monedas[i] > monedas[j - 1]:
        return solucion(opt, monedas, s, i + 1, j - 1)
    else:
        return solucion(opt, monedas, s, i, j - 2)
\end{lstlisting}

Siguendo estos pasos recursivamente obtenemos la secuencia de monedas que Sophia deberia elegir para obtener la ganancia maxima.
